\item En el análisis por activación de neutrones, una muestra es bombardeada por un flujo de neutrones constante, produciendo núcleos radiactivos a una tasa de creación constante $R$, los cuales decaen según su vida media.
\begin{enumerate}
    \item Si el bombardeo inicia a $t = 0$, demuestre que el número de átomos radiactivos acumulados al tiempo $t$ está dado por:
    $$ N(t) = \frac{R}{\lambda} \left( 1 - e^{-\lambda t} \right) $$
    \item Determine cuál es el número máximo de átomos radiactivos que es teóricamente posible acumular bajo estas condiciones.
\end{enumerate}